\documentclass[10pt,a4paper]{amsart}
\usepackage[margin=19mm]{geometry}
\usepackage{amsmath,amssymb,mathtools}
\usepackage[scr=rsfs,cal=euler]{mathalfa}
\usepackage{xfrac}
\usepackage[unicode,hidelinks,bookmarks=false]{hyperref}
\newcommand{\ie}{i.e.\ }
\newcommand{\cf}{cf.\ }
\newcommand{\resp}{respectively\ }
\newcommand{\Subsection}{\subsection}
\providecommand{\nobreakdash}{\nobreak\hbox{-}}
\setlength{\emergencystretch}{2em}
\multlinegap0pt
\usepackage{bm}
\usepackage{bbm}
\usepackage{dsfont}
\usepackage{xspace}
\usepackage{cancel}
\usepackage{mathtools}
\usepackage{amsthm}
\makeatletter
\@ifundefined{theorem}{\newtheorem{theorem}{Theorem}}{}
\makeatother
\newcommand{\C}{\mathcal{C}}
\newcommand{\F}{\mathbb{F}}
\renewcommand{\S}{\mathcal{S}}
\title[Almost Affine Vector Rank-Metric Codes]{Almost Affine Vector Rank-Metric Codes}
\author{Matteo Bonini, Johan Vester Dinesen}
\date{Source: arXiv:2605.28072v2 (CC BY 4.0)}
\begin{document}
\maketitle

\noindent\emph{Source and licence.} Almost Affine Vector Rank-Metric Codes. Version arXiv:2605.28072v2 (CC BY 4.0), distributed under the Creative Commons Attribution 4.0 International licence, as recorded in the arXiv licence metadata for this exact version and shown on the abstract page https://arxiv.org/abs/2605.28072v2. Full rights notice: licenses/arxiv-2605.28072-CC-BY-4.0.txt.\par\smallskip

\noindent\textit{Editorial scope.} Counted unit: the Theorem whose source label is \texttt{thm:constructionkdim} in the article (source0.tex lines 789--794, source label \texttt{thm:constructionkdim}), delivered together with the article's own complete original proof of it (source0.tex lines 795--825, one \texttt{proof} environment of 31 lines). No mathematics is written, repaired or re-derived: statement and proof are the article's own text line for line. Required context retained uncounted: thm:construction2dim (lines 755--757). Every definition, display and label of those lines is retained unchanged. Omitted from this package and not redistributed: 1--754, 758--788, 826--841. Nothing inside a retained excerpt is omitted or altered: every retained body is delivered exactly as the article's lines are written, so no omission receipt of any kind accompanies this package. Citations: the retained excerpts carry no citation command at all, so no bibliography entry of the article is transcribed and no reference is redistributed. Reference adaptation: none was needed and none was made; every reference inside the delivered unit resolves inside it, so no unresolved reference or citation is produced. Printed slips kept literal: no printed word, symbol, hypothesis or number of the counted statement or of its proof was altered. Layout and class: the article's own class is \texttt{article} with packages including inputenc, amsmath, amssymb, mathrsfs, mathtools, enumerate, enumitem, amsthm, booktabs, xcolor, url, hyperref, multirow, caption; the wrapper is the lane's pinned \texttt{amsart} editorial wrapper at 10pt on A4, which restates the theorem-like environments the retained text needs and adds the article's own macro definitions that the retained text needs (\texttt{\textbackslash C}, \texttt{\textbackslash F}, \texttt{\textbackslash S}), with extra wrapper packages (bm, bbm, dsfont, xspace, cancel, mathtools). The delivered numbering is the unit's own. Assets: no retained span contains a figure environment, a graphics-inclusion command, a PDF-page inclusion, a table or a TikZ picture; no image file is copied into this package, the package declares no asset dependency and no image file is redistributed with it. Licence: version arXiv:2605.28072v2 of this article is distributed under the Creative Commons Attribution 4.0 International licence, as recorded in the arXiv licence metadata for this exact version and shown on the abstract page https://arxiv.org/abs/2605.28072v2. A full rights notice is delivered at licenses/arxiv-2605.28072-CC-BY-4.0.txt. Characters: every retained excerpt is pure ASCII, so the delivered engine, XeLaTeX with its default Latin Modern text font, reports no missing character.
\begin{theorem}
Let $\S$ be a finite proper semifield of order $q^m$, and let $$\C_\S = \{ (x,(y-x\circ a))_{a\in \S}\mid (x,y)\in \S^2\}.$$ Then $\C_\S \subseteq \F_{q^m}^n$ is an $\F_q$-linear, strictly almost affine rank-metric code with parameters $n=q^m+1$, $\dim \C_\S = 2$, and $d(\C_\S)=1$. \label{thm:construction2dim}
\end{theorem}

\begin{theorem} Let $\S$ be a finite proper semifield of order $q^m$, and let $k\geq 2$. Choose a point $p =(p_1,\ldots,p_{k-1})\in \S^{k-1}$ such that there exist elements $\gamma,x\in \S$ satisfying $\gamma \circ (x\circ p_j) \neq (\gamma\circ x)\circ p_j$ for some coordinate index $j$. Define
\[
    \C_{\S,k}(p) = \left\{ \left(x_1,\ldots,x_{k-1},\left( y - \sum^{k-1}_{i=1}x_i \circ (\lambda p_i) \right)_{\lambda \in \F_q} \right) \middle| (x_1,\ldots,x_{k-1},y)\in \S^k \right\}.
\]
Then $\C_{\S,k}(p)$ is an $\F_q$-linear, strictly almost affine rank-metric code with parameters $n=(k-1)+q$, $\dim \C_{\S,k}(p) = k$, and $d(\C_{\S,k}(p)) = 1$. \label{thm:constructionkdim}
\end{theorem}

\begin{proof}
The parameters and $\F_q$-linearity follow directly by the same arguments as in Theorem~\ref{thm:construction2dim}. To prove $\mathcal{C}_{\mathcal{S},k}(p)$ is not $\mathbb{F}_{q^m}$-linear, suppose for contradiction that it is. Let $\cdot$ denote the standard field multiplication of $\F_{q^m}$. Consider the codeword $c$ generated by $(x_1, \dots, x_{k-1}, y)$. Matching the first $k-1$ coordinates and the evaluation at $\lambda = 0$ forces $\gamma \cdot c$ to equal the encoding of $(\gamma \cdot x_1, \dots, \gamma \cdot x_{k-1}, \gamma \cdot y)$. Equating the remaining coordinates therefore requires
    $$\sum_{i=1}^{k-1} \gamma \cdot (x_i \circ (\lambda p_i)) = \sum_{i=1}^{k-1} (\gamma \cdot x_i) \circ (\lambda p_i)$$
for all $\lambda \in \mathbb{F}_q$. By hypothesis, there exist an index $j$ and elements $\gamma, x \in \mathcal{S}$ such that $\gamma \circ (x \circ p_j) \neq (\gamma \circ x) \circ p_j$. Setting $x_j = x$, $x_i = 0$ for $i \neq j$, and $y = 0$, the condition at $\lambda = 1$ reduces to $\gamma \cdot (x \circ p_j) = (\gamma \cdot x) \circ p_j$. Because the proper semifield multiplication $\circ$ is strictly distinct from the field multiplication $\cdot$, this equality fails. The scaled vector is therefore not in the code, providing the contradiction.

To prove the almost affine property, consider a projection $\pi_V$ defined by an $\F_q$-subspace $V\leq \F_{q}^n$. A vector $v\in V$ is indexed as $v= (v_{\infty_1},\ldots,v_{\infty_{k-1}},(v_\lambda)_{\lambda \in \F_q})$. The condition $\langle c,v\rangle =0$ then expands to
\[
    \sum^{k-1}_{i=1} v_{\infty_i} x_i + \sum_{\lambda \in \F_q} v_\lambda \left(y - \sum^{k-1}_{i=1} x_i \circ (\lambda p_i)\right) =0.
\]
Because the scalars $\lambda \in \F_q$ associate and commute with all elements in $\S$, we have $x_i \circ (\lambda p_i) = \lambda (x_i \circ p_i)$. Regrouping terms yields
\begin{align}
    A_v y + \sum^{k-1}_{i=1} x_i \circ (v_{\infty_i} 1_\S - C_v p_i) = 0, \label{eq:rank1_kernel}
\end{align}
where $A_v = \sum_\lambda v_\lambda \in \F_q$ and $C_v = \sum_\lambda \lambda v_\lambda \in \F_q$. Define the map $C\colon V\rightarrow \F_q$ by $C(v) = C_v$. Because the coordinate summations are linear over the base field, $C$ is a well-defined linear function on $V$. Let $V_0 = \ker C$. Thus, $\dim_{\F_q} V_0 \geq \dim_{\F_q} V -1$. For any $v\in V_0$ we have $C_v =0$, so \eqref{eq:rank1_kernel} reduces to
\begin{align} \label{eq:v0_system}
    A_v y + \sum^{k-1}_{i=1} v_{\infty_i} x_i=0.
\end{align}

Let $W\subseteq \S^k$ denote the space of solutions $(x_1,\ldots,x_{k-1},y)$ of \eqref{eq:v0_system} for all $v\in V_0$. We then claim that $W$ is a left vector space over $\S$. Closure under vector addition follows immediately from the distributivity of the semifield multiplication over addition. To verify closure under left-multiplication, suppose $u=(x_1,\ldots,x_{k-1},y)\in W$, and let $s\in \S$. Substituting $s\circ u$ into the linear form then yields for any $v\in V_0$
\begin{align*}
    A_v(s\circ y) + \sum^{k-1}_{i=1}v_{\infty_i} (s\circ x_i) &= s\circ (A_v y) + \sum^{k-1}_{i=1} s\circ (v_{\infty_i} x_i) \\
    &= s\circ \left( A_v y + \sum^{k-1}_{i=1} v_{\infty_i} x_i\right) \\ &=0
\end{align*}
Thus, $W$ is a left vector space over $\S$, so its cardinality is a power of $|\S|=q^m$. 

Now, if $V = V_0$, then $\ker \pi_V = W$, and we are done. If $V \neq V_0$, the kernel $V_0$ has codimension 1 in $V$. We can therefore choose a single vector $w \in V \setminus V_0$ such that $V = V_0 \oplus \langle w \rangle_{\mathbb{F}_q}$. The total kernel of the projection by $V$ is the intersection of the $\S$-vector space $W$ with the single additional constraint imposed by $w$, which is
\[
    A_w y + \sum^{k-1}_{i=1} x_i \circ (w_{\infty_i} 1_\S - C_w p_i)=0.
\]
As $W$ is a left vector space over $\S$ and $\S$ is a division ring, adding this single linear constraint to $W$ will either be redundant and contain all of $W$, or it will uniquely determine one variable over $\S$ in terms of the others. This removes exactly one degree of freedom over $\S$, reducing the total number of solutions by a factor of $q^m$. In either case, the size of the final kernel remains a power of $q^m$.
\end{proof}

\end{document}
